% =====================================================
% TRANG GIỚI THIỆU CÁC QUYỂN TRONG BỘ SÁCH
% Truyện Tích Cảm Hứng Song Ngữ
% =====================================================
% Colors used in catalog (self-contained)
\definecolor{catred}{HTML}{8B1A1A}
\definecolor{catblue}{HTML}{1F4E79}
\definecolor{catgold}{HTML}{B8860B}
\clearpage
\thispagestyle{empty}

\begin{center}
{\large\bfseries\color{catred} CÁC QUYỂN TRONG BỘ SÁCH}\\[4pt]
{\normalsize\itshape\color{catblue} Truyện Tích Cảm Hứng Song Ngữ}\\[8pt]
{\color{catgold}\rule{10cm}{1.5pt}}
\end{center}

\vspace{6pt}

% === 9 QUYỂN ĐẦU (hiển thị đầy đủ) ===
\begin{tcolorbox}[
  colback=catred!4,
  colframe=catred!60,
  title={\bfseries Bộ Điển Tích \& Chuyện Kể (Quyển I -- IX)},
  coltitle=white,
  fonttitle=\bfseries\large,
  arc=6pt, boxrule=1pt,
  left=10pt, right=10pt, top=8pt, bottom=8pt
]
\begin{enumerate}[leftmargin=1.4em, itemsep=4pt, parsep=0pt]
  \item \textbf{Điển Tích \& Chuyện Kể Truyền Cảm Hứng} ---
    \textit{Lòng biết ơn, can đảm, tình bạn nghĩa khí và ý chí vươn lên.}

  \item \textbf{Điển Tích \& Chuyện Kể: Đạo Làm Người} ---
    \textit{Đạo làm con, làm bạn, làm trò, làm thầy và đối nhân xử thế.}

  \item \textbf{Điển Tích \& Chuyện Kể: Thiên Nhiên Quanh Ta} ---
    \textit{Bài học từ thiên nhiên và muôn loài qua câu chuyện song ngữ.}

  \item \textbf{Điển Tích \& Chuyện Kể: Triết Lý Nhân Sinh} ---
    \textit{Triết lý nhân sinh và đối nhân xử thế từ cổ chí kim.}

  \item \textbf{Điển Tích \& Chuyện Kể: Hai Mặt Của Một Đồng Xu} ---
    \textit{Tư duy biện chứng, âm dương, mâu thuẫn và cân bằng.}

  \item \textbf{Vòng Quay Cuộc Đời} ---
    \textit{Nhân quả, bản sắc, buông bỏ và người thầy vô hình.}

  \item \textbf{Lời Dạy Của Bậc Hiền Nhân} ---
    \textit{Triết lý và lời dạy từ bậc thầy qua các thời đại Đông -- Tây.}

  \item \textbf{Thô Nhưng Thật: Chuyện Học Sinh} ---
    \textit{Điểm số, áp lực, so sánh, bắt nạt --- những chuyện học đường thật.}

  \item \textbf{Ra Đời Không Có Bản Nháp} ---
    \textit{Công việc đầu tiên, tiền đầu tiên, bị từ chối, sống xa nhà.}
\end{enumerate}
\end{tcolorbox}

\vspace{8pt}

% === 9 QUYỂN SAU (hiển thị nhỏ gọn) ===
\begin{tcolorbox}[
  colback=catblue!4,
  colframe=catblue!50,
  title={\bfseries Bộ Nâng Cao \& Cuốn Tổng Hợp (Quyển X -- XVII, Cuốn 1 -- 2)},
  coltitle=white,
  fonttitle=\bfseries,
  arc=6pt, boxrule=1pt,
  left=8pt, right=8pt, top=6pt, bottom=6pt,
  fontupper=\small
]
\begin{enumerate}[leftmargin=1.4em, itemsep=2pt, parsep=0pt, start=10]
  \item \textbf{Những Ảo Tưởng Của Tuổi Trẻ} --- \textit{Nhận ra và điều chỉnh ảo tưởng trước khi đặt cược cả cuộc đời vào đó.}
  \item \textbf{Những Cái Bẫy Người Trẻ Tự Bước Vào} --- \textit{Tự mãn, cái tôi, so sánh, trì hoãn --- những bẫy ta tự giăng cho mình.}
  \item \textbf{Những Điều Thầy Giáo Nhìn Thấy} --- \textit{Bục giảng nhìn thấy những gì học sinh chưa tự thấy ở mình.}
  \item \textbf{Những Câu Hỏi Học Trò Không Nói Thẳng} --- \textit{Trả lời thẳng những câu hỏi bị giữ trong im lặng.}
  \item \textbf{100 Câu Hỏi Học Trò Hay Hỏi Về Chuyện Học} --- \textit{Chuyện học, động lực, thói quen, áp lực và con đường tương lai.}
  \item \textbf{Những Điều Thầy Muốn Nói Với Học Sinh} --- \textit{Thầy nói chuyện thẳng, gần và chậm về học tập, tính cách, đường đời.}
  \item \textbf{Những Bài Học Nhỏ Trong Lớp Học} --- \textit{Bài học lớn từ tình huống nhỏ xảy ra mỗi ngày trên lớp.}
  \item \textbf{Những Điều Học Sinh Nên Hiểu Sớm} --- \textit{Nền tảng quan trọng --- hiểu sớm thì đỡ vấp về sau.}
\end{enumerate}

\medskip
\textbf{Cuốn 1:} \textit{Những Lời Thầy Nói Trong Lớp Học} --- 50 câu lời thầy, 100 câu im lặng suy nghĩ, 50 câu hài hước sâu sắc.\\[3pt]
\textbf{Cuốn 2:} \textit{Giật Mình --- Đỉnh Cao --- Lời Cuối Năm} --- 30 câu giật mình, 50 câu đỉnh cao sư phạm, 50 câu cuối năm xúc động.
\end{tcolorbox}

\vspace{6pt}
\begin{center}
{\color{catgold}\rule{10cm}{1.5pt}}\\[4pt]
{\small\itshape\color{catblue!70} Bộ sách song ngữ Việt -- Anh dành cho giáo viên và học sinh}
\end{center}
\clearpage
\thispagestyle{empty}

\begin{center}
{\large\bfseries\color{slateink} CÁC QUYỂN TRONG BỘ SÁCH}\\[4pt]
{\normalsize\itshape\color{brickred} Truyện Tích Cảm Hứng Song Ngữ}\\[8pt]
{\color{dustgold}\rule{10cm}{1.5pt}}
\end{center}

\vspace{6pt}

% === 9 QUYỂN ĐẦU (hiển thị đầy đủ) ===
\begin{tcolorbox}[
  colback=slateink!4,
  colframe=slateink!60,
  title={\bfseries Bộ Điển Tích \& Chuyện Kể (Quyển I -- IX)},
  coltitle=white,
  fonttitle=\bfseries\large,
  arc=6pt, boxrule=1pt,
  left=10pt, right=10pt, top=8pt, bottom=8pt
]
\begin{enumerate}[leftmargin=1.4em, itemsep=4pt, parsep=0pt]
  \item \textbf{Điển Tích \& Chuyện Kể Truyền Cảm Hứng} ---
    \textit{Lòng biết ơn, can đảm, tình bạn nghĩa khí và ý chí vươn lên.}

  \item \textbf{Điển Tích \& Chuyện Kể: Đạo Làm Người} ---
    \textit{Đạo làm con, làm bạn, làm trò, làm thầy và đối nhân xử thế.}

  \item \textbf{Điển Tích \& Chuyện Kể: Thiên Nhiên Quanh Ta} ---
    \textit{Bài học từ thiên nhiên và muôn loài qua câu chuyện song ngữ.}

  \item \textbf{Điển Tích \& Chuyện Kể: Triết Lý Nhân Sinh} ---
    \textit{Triết lý nhân sinh và đối nhân xử thế từ cổ chí kim.}

  \item \textbf{Điển Tích \& Chuyện Kể: Hai Mặt Của Một Đồng Xu} ---
    \textit{Tư duy biện chứng, âm dương, mâu thuẫn và cân bằng.}

  \item \textbf{Vòng Quay Cuộc Đời} ---
    \textit{Nhân quả, bản sắc, buông bỏ và người thầy vô hình.}

  \item \textbf{Lời Dạy Của Bậc Hiền Nhân} ---
    \textit{Triết lý và lời dạy từ bậc thầy qua các thời đại Đông -- Tây.}

  \item \textbf{Thô Nhưng Thật: Chuyện Học Sinh} ---
    \textit{Điểm số, áp lực, so sánh, bắt nạt --- những chuyện học đường thật.}

  \item \textbf{Ra Đời Không Có Bản Nháp} ---
    \textit{Công việc đầu tiên, tiền đầu tiên, bị từ chối, sống xa nhà.}
\end{enumerate}
\end{tcolorbox}

\vspace{8pt}

% === 9 QUYỂN SAU (hiển thị nhỏ gọn) ===
\begin{tcolorbox}[
  colback=brickred!4,
  colframe=brickred!50,
  title={\bfseries Bộ Nâng Cao \& Cuốn Tổng Hợp (Quyển X -- XVII, Cuốn 1 -- 2)},
  coltitle=white,
  fonttitle=\bfseries,
  arc=6pt, boxrule=1pt,
  left=8pt, right=8pt, top=6pt, bottom=6pt,
  fontupper=\small
]
\begin{enumerate}[leftmargin=1.4em, itemsep=2pt, parsep=0pt, start=10]
  \item \textbf{Những Ảo Tưởng Của Tuổi Trẻ} --- \textit{Nhận ra và điều chỉnh ảo tưởng trước khi đặt cược cả cuộc đời vào đó.}
  \item \textbf{Những Cái Bẫy Người Trẻ Tự Bước Vào} --- \textit{Tự mãn, cái tôi, so sánh, trì hoãn --- những bẫy ta tự giăng cho mình.}
  \item \textbf{Những Điều Thầy Giáo Nhìn Thấy} --- \textit{Bục giảng nhìn thấy những gì học sinh chưa tự thấy ở mình.}
  \item \textbf{Những Câu Hỏi Học Trò Không Nói Thẳng} --- \textit{Trả lời thẳng những câu hỏi bị giữ trong im lặng.}
  \item \textbf{100 Câu Hỏi Học Trò Hay Hỏi Về Chuyện Học} --- \textit{Chuyện học, động lực, thói quen, áp lực và con đường tương lai.}
  \item \textbf{Những Điều Thầy Muốn Nói Với Học Sinh} --- \textit{Thầy nói chuyện thẳng, gần và chậm về học tập, tính cách, đường đời.}
  \item \textbf{Những Bài Học Nhỏ Trong Lớp Học} --- \textit{Bài học lớn từ tình huống nhỏ xảy ra mỗi ngày trên lớp.}
  \item \textbf{Những Điều Học Sinh Nên Hiểu Sớm} --- \textit{Nền tảng quan trọng --- hiểu sớm thì đỡ vấp về sau.}
\end{enumerate}

\medskip
\textbf{Cuốn 1:} \textit{Những Lời Thầy Nói Trong Lớp Học} --- 50 câu lời thầy, 100 câu im lặng suy nghĩ, 50 câu hài hước sâu sắc.\\[3pt]
\textbf{Cuốn 2:} \textit{Giật Mình --- Đỉnh Cao --- Lời Cuối Năm} --- 30 câu giật mình, 50 câu đỉnh cao sư phạm, 50 câu cuối năm xúc động.
\end{tcolorbox}

\vspace{6pt}
\begin{center}
{\color{dustgold}\rule{10cm}{1.5pt}}\\[4pt]
{\small\itshape\color{brickred!70} Bộ sách song ngữ Việt -- Anh dành cho giáo viên và học sinh}
\end{center}
\clearpage
